\begin{abstract}
The rapid advancement of generative artificial intelligence has catalyzed a paradigm shift in quantitative finance, transitioning from traditional point forecasting toward comprehensive financial time-series generation (FTSG) across extrapolation, imputation, and full-scale data synthesis. This transition is motivated by the critical need to simulate rare market regimes, reconstruct incomplete historical records, and produce privacy-preserving synthetic environments for stress testing and trading execution. However, empirical financial properties---including heavy tails, volatility clustering, non-stationarity, and regime shifts---pose severe hurdles for generic generative frameworks. In this systematic survey, we review \revdeldec{ 130 }\revadddec{131 }studies from 2021 to early 2026 to bridge the gap between machine learning capabilities and econometric validity. We introduce a two-level taxonomy structuring the domain across task formulations and algorithmic families, and establish a multi-dimensional evaluation protocol encompassing distributional fidelity, stylized facts preservation, and downstream utility. Furthermore, we catalog curated financial databases and benchmark suites while addressing empirical dataset construction pitfalls. Finally, we identify pivotal open challenges in causal representation learning, econometrically bounded modeling, and online adaptation under distribution shifts. This review provides a standardized roadmap for developing robust, domain-aligned financial generative models.
\end{abstract}